\section{Density and small connected districts}
\label{sec:density}

The sharp order of growth of the density threshold for clique minors
was established by Kostochka \cite{Kostochka1982,Kostochka1984} and
Thomason \cite{Thomason1984}. We prove below a coarser estimate sufficient for
our purpose. We need two consequences of density. The incidence estimate of
Lemma~\ref{lem:incidence} will control exceptional vertices in the
path-history bound, Lemma~\ref{separate:footprint}, after paths are
contracted. Lemma~\ref{lem:district} supplies small highly connected
districts in which the later constructions can route paths and build
clique models.

Write \(e(R)=|E(R)|\) and \(d(R)=e(R)/|V(R)|\) for a nonempty graph \(R\).
Throughout this section, fix
\[
 \epsilon=\frac1{200}.
\]
All assertions involving sufficiently large \(t\) have a cutoff independent of
the graph. When a fixed parameter \(K\) is present, the cutoff may depend on
\(K\). An upper bound on an integer order by a real power of \(t\) is understood
as an ordinary real inequality.

\begin{lemma}\label{lem:coarse-density}
There is an absolute constant \(C_*>0\), for example \(C_*=10^4\),
such that, for every integer \(t\ge2\)
and every nonempty graph \(P\) with \(h(P)<t\),
\[
 d(P)\le C_*t\log(2t).
\]
Consequently, every graph \(P\) with \(h(P)<t\) satisfies
\[
 \ell(P)\le 2C_*t\log(2t)+1.
\]
\end{lemma}

\begin{proof}
Choose a nonempty minor \(R\) of \(P\) whose density \(d\) is as large as
possible. If \(d<10^4\), the claimed density estimate holds with
\(C_*=10^4\). We may therefore assume \(d\ge10^4\).
Put \(n=|V(R)|\). The graph \(R\) has no isolated vertex, since deleting
one would increase its positive density.
Contracting an edge \(xy\) with \(c_{xy}\) common neighbors
loses exactly \(1+c_{xy}\) edges. The maximality of \(d\) gives
\[
 \frac{dn-1-c_{xy}}{n-1}\le d,
 \qquad\text{and hence}\qquad c_{xy}\ge d-1.
\]
There is a vertex \(v\) of degree at most \(2d\). Therefore
\(J=R[N_R(v)]\) has at most \(2d\) vertices and minimum degree at least \(d-1\).

We first record a bound for connected dominating sets in any connected graph
\(C\) of order at most \(2d\) and minimum degree at least \(d/2\).
Take
\[
 m=\lceil4\log(2d)\rceil+1
\]
independent uniform samples from \(V(C)\). A fixed vertex has a neighbor in
each draw with probability at least \(1/4\), so the probability that some
vertex has no sampled neighbor is at most
\[
 2d(3/4)^m\le 2d\exp(-m/4)<1.
\]
Consequently, some choice of at most \(m\) sampled vertices dominates \(C\).
The diameter of \(C\) is at most \(11\): along a shortest path, the open
neighborhoods of the vertices at positions \(0,3,6,\ldots\) are pairwise
disjoint, each has at least \(d/2\) vertices, and the graph has at most \(2d\)
vertices. There are at most four such positions. Joining each sampled
vertex to one fixed sampled vertex along a shortest path therefore gives
a connected dominating set of size at most
\(11m\le66\log(2d)\), which is at most
\[
 M=\lceil100\log(2d)\rceil.
\]
We have \(M\le101\log(2d)\) and \(d\ge6M\).
For the latter inequality, \(x/\log(2x)\) is increasing for \(x\ge10^4\),
and \(10^4>606\log(20000)\).

We now construct
\[
 q=\left\lfloor\frac{d}{3M}\right\rfloor
\]
pairwise adjacent, disjoint connected sets. Start with \(J_0=J\).
At step \(i\), choose a component \(C_i\) of \(J_i\), take a connected
dominating set \(B_i\) of \(C_i\) of size at most \(M\), and put
\(J_{i+1}=C_i-B_i\).
For every retained vertex, passing to a component deletes no neighbor.
Its neighbors in the original graph \(J\) that have disappeared are
therefore all in previously chosen bags. As long as at most \(qM\le d/3\)
bag vertices have been deleted, each retained vertex has degree at least
\[
 d-1-qM\ge d/2.
\]
In particular, each chosen component has at least \(d/2+1\) vertices,
whereas \(M\le d/6\); its deletion of \(B_i\) leaves a nonempty graph.
The construction consequently continues for all \(q\) steps.

Every later bag lies inside \(C_i-B_i\), so it has an edge to \(B_i\)
by domination. The bags form a \(K_q\) model. Since \(d/(3M)\ge2\),
\[
 h(P)\ge q\ge\frac{d}{6M}
 \ge\frac{d}{606\log(2d)}.
\]
The hypothesis \(h(P)<t\) implies \(d/\log(2d)<606t\).
The elementary inequality \(\log(2d)\le2\sqrt d\) for \(d\ge1\) first
gives \(d<(1212t)^2\). Since \(2\cdot1212^2<2^{22}\) and \(t\ge2\),
\[
 \log(2d)<22\log2+2\log t\le12\log(2t).
\]
It follows that \(d<7272t\log(2t)\). Together with the case \(d<10^4\)
and the inequality \(d(P)\le d\), this proves the density claim with
\(C_*=10^4\).

Every nonempty subgraph of \(P\) has average degree at most
\(2C_*t\log(2t)\), and hence has a vertex of degree at most
\(\lfloor2C_*t\log(2t)\rfloor\).
Successively deleting such vertices and coloring in the reverse order
proves the asserted list-coloring bound. The empty graph causes no exception.
\end{proof}

\begin{lemma}\label{lem:incidence}
There is an absolute positive integer \(B_*\) such that, for all sufficiently
large integers \(t\), every graph \(P\) with \(h(P)<t\) and disjoint vertex
sets \(A,T\) satisfies
\begin{equation}\label{density:incidence-bound}
 \sum_{\substack{v\in T\\ |N_P(v)\cap A|\ge B_*t}}
       |N_P(v)\cap A|
 \le t^{1+2\epsilon}|A|.
\end{equation}
In particular,
\[
 \bigl|\{v\in T:|N_P(v)\cap A|\ge B_*t\}\bigr|
 \le t^{2\epsilon}|A|.
\]
\end{lemma}

\begin{proof}
Choose a fixed positive even integer \(s\) with \(\epsilon s/2>4\), and
put \(B_*=2s\). We first show that, in any graph without a \(K_t\) minor,
every set \(U\) of size \(b\ge B_*t\) contains a vertex with at least
\(t^{-\epsilon}b\) nonneighbors in \(U\) distinct from itself.
Suppose otherwise, and set \(\delta=t^{-\epsilon}\).
Sample an ordered sequence of \(st\) distinct vertices of \(U\) uniformly,
and divide its positions into \(t\) batches of size \(s\).

After any previously exposed sampled positions have been fixed, a further
position is uniform among at least \(b-st\ge b/2\) remaining vertices.
For any already exposed vertex, the conditional probability that this
new vertex is a nonneighbor is consequently at most \(2\delta\).
Consider a fixed nontrivial cut of one batch. Expose a fixed position on
its smaller side first, and then the positions on its larger side.
If the cut has no crossing edge, at least \(s/2\) successive exposures
must all be nonneighbors of the first vertex. The probability is at most
\((2\delta)^{s/2}\). Taking the union over at most \(2^s\) cuts, the probability
that the batch is disconnected is at most \(2^s(2\delta)^{s/2}\).
For two fixed batches to be anticomplete, every position of the second
batch must be a nonneighbor of one fixed position of the first; this
has probability at most \((2\delta)^s\).
Thus the probability that some batch is disconnected or some two batches
are anticomplete is at most
\[
 t2^s(2t^{-\epsilon})^{s/2}
 +\binom t2(2t^{-\epsilon})^s=o(1).
\]
For large \(t\), the batches can therefore be chosen connected and pairwise
adjacent, giving a \(K_t\) model. This contradiction proves the claim.

Let \(T_0\) consist of the vertices appearing in the sum in
\eqref{density:incidence-bound}. If \(A\) or \(T_0\) is empty, there is
nothing to prove. Otherwise, delete all vertices outside \(A\cup T_0\)
and all edges with both ends in \(T_0\).
Process the vertices of \(T_0\) one at a time. Throughout the procedure,
the vertices representing \(A\) carry a graph that only gains edges.
Each unprocessed vertex retains precisely its original neighborhood in
\(A\), because there were no edges between vertices of \(T_0\).

For an unprocessed vertex \(v\), put \(b_v=|N_P(v)\cap A|\).
Apply the preceding claim to its neighborhood in the current graph on
\(A\), which is a subgraph of a minor of \(P\). Choose there a vertex
with at least \(t^{-\epsilon}b_v\) missing neighbors in that neighborhood,
and contract \(v\) onto this vertex. This operation adds at least
\(t^{-\epsilon}b_v\) new edges to the graph on \(A\).
After all vertices have been processed, the resulting graph on \(A\)
is a minor of \(P\). By Lemma~\ref{lem:coarse-density},
\[
 t^{-\epsilon}\sum_{v\in T_0}b_v
 \le e_{\mathrm{final}}-e(P[A])
 \le C_*t\log(2t)|A|.
\]
For sufficiently large \(t\), \(C_*\log(2t)\le t^\epsilon\), proving
\eqref{density:incidence-bound}. Dividing by \(B_*t\), and then using
\(B_*\ge1\), proves the last assertion.
\end{proof}

The following density-to-connectivity lemma is Mader's theorem
\cite{Mader1972}; see also \cite[Lemma~4.4]{DelcourtPostle2021}.
We include the elementary proof needed here.

\begin{lemma}\label{density:connectivity}
For every positive integer \(k\), a graph of density greater than \(2k\)
contains an induced \(k\)-connected subgraph.
\end{lemma}

\begin{proof}
Choose an induced subgraph \(H\) of smallest order \(v\ge2k\) such that
\[
 e(H)>2k(v-k).
\]
Such a subgraph exists: the given graph has density greater than \(2k\)
and, being simple, has more than \(4k+1\) vertices.
The equality \(v=2k\) is impossible because
\(\binom{2k}{2}<2k^2\).
Deleting any one vertex therefore leaves at least \(2k\) vertices,
and minimality implies
\[
 \deg_H(x)
 =e(H)-e(H-x)
 >2k(v-k)-2k(v-1-k)=2k
\]
for every \(x\in V(H)\).

Suppose that \(H\) is not \(k\)-connected. There is a separator \(Z\)
of size \(z<k\) and nonempty anticomplete sides \(A,B\) partitioning
\(V(H)\setminus Z\).
Every neighbor of a vertex of \(A\) lies in \(A\cup Z\), and likewise
for \(B\). Both \(H[A\cup Z]\) and \(H[B\cup Z]\) therefore have at
least \(2k\) vertices and are proper induced subgraphs of \(H\).
Their minimality bounds give
\[
 \begin{aligned}
 e(H)
 &=e(H[A\cup Z])+e(H[B\cup Z])-e(H[Z])\\
 &\le2k\bigl(|A|+z-k+|B|+z-k\bigr)-e(H[Z])\\
 &\le2k(v+z-2k)
 <2k(v-k),
 \end{aligned}
\]
a contradiction. Thus \(H\) is \(k\)-connected.
\end{proof}

The next argument uses the maximal-contraction method of Delcourt and
Postle~\cite{DelcourtPostle2021}. The edge-loss estimates below give the
particular order and connectivity bounds needed here.

\begin{lemma}\label{lem:district}
For every fixed positive integer \(K\), there is an integer
\[
 D_K>\max(400B_*,12800K)
\]
such that, for all sufficiently large integers \(t\), every nonempty
graph \(P\) with
\[
 h(P)<t,\qquad d(P)\ge D_Kt
\]
contains an induced \((Kt)\)-connected subgraph of order at most
\(t^{1.06}\).
\end{lemma}

\begin{proof}
Choose an integer \(D_K\) satisfying the displayed inequality, and keep
it fixed. Write \(n=|V(P)|\). Delete edges to obtain a spanning subgraph
\(P_0\) with exactly \(dn\) edges, where \(d=D_Kt\).
This is possible because \(D_Ktn\) is an integer.
Restoring deleted edges at the end preserves the connectivity of any
induced subgraph. Put
\[
 p_0=\lfloor t^{5\epsilon}\rfloor,\qquad
 \lambda=\frac d{100}.
\]
We take \(t\) large enough that \(p_0\ge2\).

We will find a small dense subgraph on vertices untouched by the
contractions. First we contract a maximum family of connected
\(p_0\)-vertex sets within an edge-loss budget. We then retain
untouched vertices whose degrees have not fallen too far. A second
maximal choice of a connected set will stop before reaching \(p_0\),
since reaching that order would enlarge the first family. The failure
of its one-vertex extensions will force a dense neighborhood, even
after the contracted vertices are removed.

Choose a family \(\mathcal C\) of as many pairwise disjoint connected
vertex sets of \(P_0\), each of size \(p_0\), as possible subject to
\begin{equation}\label{density:family-budget}
 e(P_0)-e(P_0/\mathcal C)
 \le \lambda\sum_{C\in\mathcal C}(|C|-1).
\end{equation}
Here \(P_0/\mathcal C\) denotes the simple graph obtained by contracting
each set in the family and suppressing loops and parallel edges.
The empty family is allowed, so a maximum exists.
Let \(R=P_0/\mathcal C\), and let \(X\) be the set of contracted vertices.
The vertices outside \(X\) are identified with their original vertices
in \(P_0\), and edges between two such vertices are unchanged. We have
\begin{equation}\label{density:quotient-bounds}
 |X|\le\frac n{p_0},
 \qquad
 (d-\lambda)n\le e(R)\le dn.
\end{equation}
Indeed the total vertex loss in the contractions is at most \(n\),
so \eqref{density:family-budget} gives the lower edge bound.

Let \(B\) be the union of the following three sets:
\[
 X,\qquad
 \{v\in V(R)\setminus X:|N_R(v)\cap X|\ge\lambda/4\},
 \qquad
 \{v\in V(R):\deg_R(v)\ge dt^{4\epsilon}\}.
\]
The first exclusion ensures that the vertices retained below are
original vertices. The second limits how many neighbors a retained
vertex can lose when \(X\) is removed from the eventual neighborhood.
The third bounds the order of that neighborhood by bounding the degrees
of the vertices that define it.
Since \(\lambda/4=D_Kt/400\ge B_*t\), Lemma~\ref{lem:incidence} bounds
the second set by \(t^{2\epsilon}|X|\).
The degree sum and \eqref{density:quotient-bounds} bound the third
set by \(2n/t^{4\epsilon}\).
For sufficiently large \(t\), \(p_0\ge t^{5\epsilon}/2\), and hence
\begin{equation}\label{density:bad-size}
 |B|\le(1+t^{2\epsilon})\frac n{p_0}+\frac{2n}{t^{4\epsilon}}
 \le6nt^{-3\epsilon}.
\end{equation}

By Lemma~\ref{lem:coarse-density}, the number of edges within \(B\) is
at most
\[
 C_*t\log(2t)|B|
 \le6C_*nt^{1-3\epsilon}\log(2t).
\]
To bound the edges from \(B\) to its complement, first count the
vertices outside \(B\) with fewer than \(B_*t\) neighbors in \(B\):
their contribution is at most \(B_*tn\).
Lemma~\ref{lem:incidence} bounds all remaining crossing edges by
\[
 t^{1+2\epsilon}|B|\le6nt^{1-\epsilon}.
\]
The same bounds hold if \(B\) is empty. It follows that
\[
 \begin{aligned}
 e(R-B)
 &\ge (d-\lambda)n
       -6C_*nt^{1-3\epsilon}\log(2t)
       -B_*tn-6nt^{1-\epsilon}\\
 &\ge dn\left(
 \frac{99}{100}-\frac{B_*}{D_K}
 -\frac{6C_*}{D_K}t^{-3\epsilon}\log(2t)
 -\frac6{D_K}t^{-\epsilon}\right)
 >\frac{dn}{2}
 \end{aligned}
\]
for sufficiently large \(t\), since \(D_K>400B_*\).

Starting with \(R-B\), repeatedly delete a vertex \(v\) whose degree
in the current graph is less than \(d/4\), or less than
\(\deg_R(v)/8\). Charge each deleted edge to the endpoint deleted first.
The total charge is at most
\[
 \frac d4n+\frac18\sum_{v\in V(R)}\deg_R(v)
 \le\frac{dn}{2}.
\]
The process cannot delete every vertex, because \(R-B\) initially
has more than \(dn/2\) edges. The surviving nonempty set \(S\) satisfies
\begin{equation}\label{density:surviving-degrees}
 \deg_{R[S]}(v)\ge d/4,
 \qquad
 \deg_{R[S]}(v)\ge\deg_R(v)/8
 \quad(v\in S).
\end{equation}
The second degree bound will let us compare a neighborhood in \(R\)
with the part remaining in \(S\), once the loss caused by contracting
its defining set is controlled.

For a nonempty connected vertex set \(U\) in \(R\), let
\(\mathcal L(U)\) be the number of edges lost when \(U\) is contracted
to one vertex. If \(N_R(U)\) denotes its external neighborhood, then
the exact loss is
\begin{equation}\label{density:loss-identity}
 \mathcal L(U)
 =e(R[U])+
 \sum_{w\in N_R(U)}\bigl(|N_R(w)\cap U|-1\bigr).
\end{equation}
The first term counts the internal edges that become loops.
At each outside neighbor, all but one of its edges into \(U\)
become parallel edges and are suppressed, giving the second term.

Choose a connected set \(H\subseteq S\) of largest possible order
\(u\le p_0\) subject to
\begin{equation}\label{density:new-budget}
 \mathcal L(H)\le\lambda(u-1).
\end{equation}
A singleton is admissible. We have \(u<p_0\).
Indeed \(H\) contains only untouched original vertices, so it is
connected in \(P_0\) as well. If \(u=p_0\), append it to
\(\mathcal C\). Contracting the old family and then \(H\) gives
the same graph as contracting the enlarged family.
The actual edge losses add, as do the vertex losses; combining
\eqref{density:family-budget} and \eqref{density:new-budget}
would make this a larger admissible family, a contradiction.

Write \(N=N_R(H)\), and put
\[
 a=\sum_{v\in H}\deg_{R[S]}(v),
 \qquad
 b=\sum_{v\in H}\deg_R(v).
\]
We first show that at least one sixteenth of \(N\) lies in \(S\).
By \eqref{density:loss-identity},
\[
 \begin{aligned}
 a-|N\cap S|
 &=2e(R[H])
   +\sum_{w\in N\cap S}\bigl(|N_R(w)\cap H|-1\bigr)\\
 &\le2\mathcal L(H)\le2\lambda(u-1)\le2\lambda u.
 \end{aligned}
\]
By \eqref{density:surviving-degrees},
\(a\ge du/4=25\lambda u\) and \(a\ge b/8\).
As \(|N|\le b\), this gives
\begin{equation}\label{density:neighborhood-ratio}
 |N\cap S|
 \ge a-2\lambda u
 \ge\frac{23}{25}a
 \ge\frac{23}{200}b
 \ge\frac{|N|}{16}.
\end{equation}
In particular \(N\cap S\) is nonempty, since \(a\ge25\lambda u>0\).

Fix \(w\in N\cap S\), and contract \(H\) to a vertex \(h\).
Contracting the edge \(hw\) next loses exactly
\[
 1+|N_R(w)\cap N|
\]
edges: the common neighbors of \(h\) and \(w\) are precisely
\(N_R(w)\cap N\).
If this loss were at most \(\lambda\), the connected set
\(H\cup\{w\}\), whose order is at most \(p_0\), would satisfy
the budget \(\mathcal L(H\cup\{w\})\le\lambda u\).
This would contradict the maximality of \(u\).
Thus \(w\) has more than \(\lambda-1\) neighbors in \(N\).
Since \(w\in S\) is not in \(B\), it has fewer than \(\lambda/4\)
neighbors in \(X\). For \(\lambda\ge4\), it follows that
\[
 \deg_{R[N\setminus X]}(w)
 >\frac{3\lambda}{4}-1
 \ge\frac{\lambda}{2}.
\]
Using \(N\cap S\subseteq N\setminus X\) and
\eqref{density:neighborhood-ratio}, we obtain
\[
 2e(R[N\setminus X])
 \ge\frac{\lambda}{2}|N\cap S|
 \ge\frac{\lambda}{32}|N|.
\]
Consequently,
\begin{equation}\label{density:small-neighborhood-density}
 d(R[N\setminus X])
 \ge\frac{\lambda|N|}{64|N\setminus X|}
 \ge\frac{\lambda}{64}.
\end{equation}

Every vertex of \(H\subseteq S\) has degree less than \(dt^{4\epsilon}\)
in \(R\). Therefore
\[
 |N\setminus X|
 \le |N|
 \le\sum_{v\in H}\deg_R(v)
 <u\,dt^{4\epsilon}
 \le D_Kt^{1+9\epsilon}
 =D_Kt^{1.045}.
\]
Since \(D_K\) is fixed, this is at most \(t^{1.06}\) once
\(t^{3\epsilon}\ge D_K\).
Also,
\[
 \frac{\lambda}{64}=\frac{D_Kt}{6400}>2Kt.
\]
Lemma~\ref{density:connectivity}, applied to \(R[N\setminus X]\),
now gives an induced \((Kt)\)-connected subgraph of order at most
\(t^{1.06}\). All vertices in \(N\setminus X\) are untouched original
vertices, and the edges between them in \(R\) are exactly their edges
in \(P_0\). Restoring the deleted edges on the same vertex set gives
the required induced subgraph of \(P\).
\end{proof}

\begin{corollary}\label{density:list-district}
Fix \(K\) and \(D_K\) as in Lemma~\ref{lem:district}.
For all sufficiently large integers \(t\), if
\[
 h(P)<t,\qquad \ell(P)>2D_Kt+2,
\]
then \(P\) contains an induced \((Kt)\)-connected subgraph of order
at most \(t^{1.06}\).
\end{corollary}

\begin{proof}
If every nonempty subgraph had a vertex of degree at most \(2D_Kt+1\),
reverse deletion order would color arbitrary lists of size \(2D_Kt+2\).
Thus some subgraph has minimum degree at least \(2D_Kt+2\).
Passing to the induced subgraph on its vertices preserves this degree
bound, so that induced subgraph has density at least \(D_Kt+1\).
Apply Lemma~\ref{lem:district} there.
\end{proof}
